% sec-03-general-investigational-plan.tex - General Investigational Plan (full IND)
% Phase 1 PDAC IND: AI Generation. IND v1.0; repository v4.3.0.
\section{General Investigational Plan}

This General Investigational Plan describes, at the level of detail required by
21 CFR \S 312.23(a)(3)(iv), the rationale for the investigation, the indication to
be studied, the general approach to evaluating the investigational articles, the
trials to be conducted during the first year, the estimated number of subjects to
be exposed, and the foreseeable drug-related and device-related risks together with
their mitigations. The investigation is a single, integrated, first-in-human study
of two regulated articles administered in combination: perioperative oral
daraxonrasib (RMC-6236), an investigational pan-RAS(ON) multi-selective
small-molecule inhibitor evaluated under the investigational new drug pathway, and
an on-premises large language model (LLM) directed eight-arm robotic
pancreaticoduodenectomy (Whipple) system, a significant-risk Class II collaborative
device evaluated under the investigational device exemption pathway. The study
carries a single protocol, a single informed consent bearing the Physical AI
opt-out, and a single three-tier safety governance structure. The plan that follows
is intentionally quantitative so that the reviewer can assess the design, the dose
schedule, the operating characteristics, and the benefit-risk balance from the
numbers themselves \cite{phase1guidance,daraxwhipple}.

\subsection{Rationale}

The scientific rationale rests on a population of profound unmet need and a
mechanistically matched intervention. Pancreatic ductal adenocarcinoma (PDAC)
remains among the most lethal of solid tumors: the five-year overall survival is
below 13 percent, and the disease is the third leading cause of cancer death in the
United States and the fourth in the European Union \cite{pdac060s2030}. The great
majority of PDAC is driven by an activating KRAS mutation at codon 12, most commonly
KRAS G12D, G12V, or G12R, and until recently this oncogenic dependency was
considered pharmacologically intractable. Daraxonrasib (RMC-6236) is an oral
once-daily RAS(ON) multi-selective (pan-RAS) inhibitor that binds the GTP-bound,
active conformation of mutant RAS in a tri-complex with cyclophilin A and thereby
suppresses downstream effector signaling through the MAPK and PI3K pathways. The
agent received FDA Breakthrough Therapy designation in June 2025 for previously
treated metastatic PDAC harboring a KRAS G12 mutation, and the 300 mg once-daily
dose is the established recommended Phase 2 dose of the RASolute 302 pan-KRAS
program \cite{daraxwhipple,oncotrialllm}. Pairing a genotype-matched systemic agent
with a margin-negative curative-intent resection in the resectable and
borderline-resectable setting is a scientifically coherent strategy for a disease in
which the resection margin is the strongest modifiable determinant of survival.

The second component of the rationale is the surgical platform itself. The Whipple
operation is the most technically demanding curative-intent abdominal procedure, and
its early morbidity is dominated by pancreatic fistula, intra-abdominal hemorrhage,
sepsis, and failure to rescue. The investigational eight-arm robotic system is
engineered to perform this operation within hard cyber-physical limits under
continuous human oversight: eight arms each with seven degrees of freedom (56
degrees of freedom total), 640 sensor channels (80 per arm) sampled at 100 kHz for
force and 10 kHz for other modalities, a positional accuracy of 0.05 mm RMS at
1200 mm/s, per-arm tip-force caps at or below 3 N and a cumulative cross-arm cap at
or below 18 N, a 10 kHz safety heartbeat with a 100 microsecond watchdog, arm-park
within 50 microseconds, a cross-arm emergency stop within 3 ms to 0.0 N, and a
hardware emergency stop within 500 ms consistent with 21 CFR \S 312.404
\cite{onpremwhippl,daraxwhipple}. A vascular no-fly gate continuously protects five
named vessels (the superior mesenteric vein, portal vein, hepatic artery, celiac
axis, and superior mesenteric artery) with a 3.0 mm advisory soft-warning zone, a
1.0 mm no-fly zone for the superior mesenteric and portal veins and a 1.5 mm no-fly
zone for the arterial structures, and a 5.0 mm hard-stop zone that triggers an
emergency stop. The on-premises LLM operates strictly as a second-opinion advisory
oracle, hash-pinned to a deterministic seed and a repository commit, isolated
outside the vendor kinematic stack, and prohibited from full autonomy consistent
with 21 CFR \S 312.21(e); every advisory passes a ten-gate
verification-before-generation suite returning ACCEPT, BLOCK, or ESCALATE, aligned
with H.R. 9510, the Verification Before Generation in Physical AI Oncology Trials
Act of 2026 \cite{congress9510,clintrustpai}.

The schedule rationale is distinct from, and complementary to, the clinical
rationale. This IND is itself the highest-value product of a verified single-prompt
authoring method whose purpose is to compress the administrative and preparation
time that separates clinical milestones, without ever altering tumor biology or any
fixed regulatory clock. Figure 1 ranks the six document targets where faster,
validated authoring yields the greatest schedule value. The initial IND and IRB
package, the subject of this build, is the highest-value target because it sits on
the critical path; one input, faster validated authoring, feeds all six targets, and
each contributes to a common compression of administrative and preparation time
measured in months to a year. Crucially, the gains are capped by clinical events and
by fixed review clocks: faster authoring cannot shorten the 30-day FDA review or
manufacture missing toxicology or stability data \cite{paipredict4x,natlpaiplatf}.

\begin{mermaidfig}
\node[mmin]    (in) at (0,-4.5)     {Faster, validated\\IND authoring};
\node[mmgoal]  (t1) at (6.0,0)      {1. Initial IND + IRB\\package (this build)};
\node[mmgoal]  (t2) at (6.0,-1.8)   {2. Protocol amendments\\+ synchronized consent};
\node[mmproc]  (t3) at (6.0,-3.6)   {3. Cohort-review package\\after DLT data mature};
\node[mmgoal]  (t4) at (6.0,-5.4)   {4. Complete clinical-hold\\response};
\node[mmaccent](t5) at (6.0,-7.2)   {5. EOP2 briefing +\\Phase 3 protocol};
\node[mmaccent](t6) at (6.0,-9.0)   {6. CSR + NDA/BLA\\modules at lock};
\node[mmgoal]  (out)at (12.0,-3.0)  {Compressed admin /\\prep time\\(months to a year)};
\node[mmdec]   (lim)at (12.0,-7.4)  {Capped by clinical\\events and fixed\\review clocks};
\draw[mmedge] (in.east) to[out=70,in=180,looseness=0.5]  (t1.west);
\draw[mmedge] (in.east) to[out=45,in=180,looseness=0.5]  (t2.west);
\draw[mmedge] (in.east) to[out=20,in=180,looseness=0.5]  (t3.west);
\draw[mmedge] (in.east) to[out=-20,in=180,looseness=0.5] (t4.west);
\draw[mmedge] (in.east) to[out=-45,in=180,looseness=0.5] (t5.west);
\draw[mmedge] (in.east) to[out=-70,in=180,looseness=0.5] (t6.west);
\draw[mmedge] (t1.east) to[out=0,in=120,looseness=0.6] (out.west);
\draw[mmedge] (t2.east) to[out=0,in=150,looseness=0.6] (out.west);
\draw[mmedge] (t3.east) -- (out.west);
\draw[mmedge] (t4.east) to[out=0,in=200,looseness=0.6] (out.west);
\draw[mmedge] (t5.east) to[out=0,in=225,looseness=0.6] (out.west);
\draw[mmedge] (t6.east) to[out=0,in=250,looseness=0.6] (out.west);
\draw[mmedged] (out.south) to[out=-90,in=90] (lim.north);
\end{mermaidfig}
\figcaption{Figure 1. The six greatest-acceleration IND document targets, ranked,
with the initial IND and IRB package the highest-value target and the subject of
this build. One input, faster validated authoring, feeds all six; each contributes
to a common compression of administrative and preparation time of months to a year.
The schedule value is highest on the critical path, and the gains are capped by
clinical events and fixed review clocks (they cannot shorten the 30-day FDA review
or manufacture missing data).}

The mechanism of the schedule benefit is specific and is shown in Figure 2. A
clinical-trial timeline divides into three buckets: a clinical and operational
bucket (recruitment, surgery, the dose-limiting-toxicity window, and
progression-free and overall-survival maturation), which is not compressible; an
administrative and preparation bucket (data cleaning, analysis, and document
authoring), which is compressible; and a fixed regulatory-review bucket (the 30-day
IND review and the 30-day clinical-hold review), which is fixed. The repository LLM
acts on the middle bucket only. The exact per-document savings are: the initial IND
and IRB package falls from 8 to 12 weeks to 1 to 4 days; a protocol amendment with
synchronized consent falls from 2 to 4 weeks to 1 to 2 days; a cohort-review package
falls from 1 to 2 weeks to about 1 day; a complete clinical-hold response falls from
3 to 6 weeks to 2 to 4 days; an annual development safety update report (DSUR) falls
from 4 to 8 weeks to 2 to 3 days; and a clinical study report under ICH E3 falls
from 3 to 6 months to 1 to 2 weeks, summing to months to a year saved cumulatively
\cite{paipredict4x,paisitedocpk}.

\begin{mermaidfig}
\node[mmdec]  (b1) at (0,0)        {Clinical /\\operational\\NOT compressible};
\node[mmaccent](b2)at (5.6,0)      {Administrative /\\preparation\\COMPRESSIBLE};
\node[mmdec]  (b3) at (11.2,0)     {Regulatory review\\FIXED};
\node[mmctx]  (b1d)at (0,-2.4)     {Recruitment, surgery,\\DLT window, PFS / OS\\maturation};
\node[mmctx]  (b3d)at (11.2,-2.4)  {30-day IND review,\\30-day clinical-hold\\review};
\node[mmproc] (llm)at (5.6,-2.4)   {Repository LLM acts\\on this bucket only};
\node[mmin]   (s1) at (1.4,-5.0)   {IND + IRB package:\\8-12 weeks to 1-4 days};
\node[mmin]   (s2) at (5.6,-5.0)   {Protocol amendment +\\consent: 2-4 weeks to\\1-2 days};
\node[mmin]   (s3) at (9.8,-5.0)   {Cohort-review package:\\1-2 weeks to ~1 day};
\node[mmin]   (s4) at (1.4,-7.4)   {Clinical-hold response:\\3-6 weeks to 2-4 days};
\node[mmin]   (s5) at (5.6,-7.4)   {Annual DSUR:\\4-8 weeks to 2-3 days};
\node[mmin]   (s6) at (9.8,-7.4)   {CSR (ICH E3):\\3-6 months to 1-2 weeks};
\node[mmgoal] (g)  at (5.6,-9.6)   {Months to a year saved\\cumulatively};
\draw[mmedge] (b1) -- (b1d);
\draw[mmedge] (b3) -- (b3d);
\draw[mmedge] (b2) -- (llm);
\draw[mmedge] (llm.south) to[out=-120,in=90,looseness=0.6] (s1.north);
\draw[mmedge] (llm.south) -- (s2.north);
\draw[mmedge] (llm.south) to[out=-60,in=90,looseness=0.6] (s3.north);
\draw[mmedge] (s1) -- (s4);
\draw[mmedge] (s2) -- (s5);
\draw[mmedge] (s3) -- (s6);
\draw[mmedge] (s4.south) to[out=-90,in=150,looseness=0.6] (g.west);
\draw[mmedge] (s5.south) -- (g.north);
\draw[mmedge] (s6.south) to[out=-90,in=30,looseness=0.6] (g.east);
\end{mermaidfig}
\figcaption{Figure 2. The three timeline buckets. Only the administrative and
preparation bucket compresses; the clinical and operational bucket and the fixed
regulatory review clocks do not. The repository LLM acts on the middle bucket only,
with exact per-document savings: initial IND and IRB package 8 to 12 weeks to 1 to 4
days; protocol amendment with synchronized consent 2 to 4 weeks to 1 to 2 days;
cohort-review package 1 to 2 weeks to about 1 day; complete clinical-hold response 3
to 6 weeks to 2 to 4 days; annual DSUR 4 to 8 weeks to 2 to 3 days; clinical study
report 3 to 6 months to 1 to 2 weeks, summing to months to a year saved
cumulatively.}

These two rationales meet at the patient. Figure 3 traces the patient time-saved
cascade: faster validated authoring advances each step from IND and IRB submission
and site activation, through the elapse of the 30-day clock and the first dose, to
faster 3+3 cohort turnover toward the recommended Phase 2 dose, shortening the white
space between phases and, for this low-survival population, extending
progression-free and overall survival, without changing tumor biology or any fixed
review clock. The cascade is the patient-level expression of the schedule benefit
and is the bridge between the rationale here and the favorable benefit-risk balance
developed in \S 4.6 \cite{paipredict4x,pdac060s2030}.

\begin{mermaidfig}
\node[mmaccent](a) at (0,0)     {Faster validated\\IND authoring};
\node[mmproc] (b) at (4.2,0)    {Earlier IND / IRB\\submission +\\activation};
\node[mmproc] (c) at (8.4,0)    {Earlier 30-day clock\\elapses; first dose};
\node[mmproc] (d) at (12.6,0)   {Faster 3+3 cohort\\turnover to RP2D};
\node[mmproc] (e) at (12.6,-2.6){Shorter white space\\between phases};
\node[mmgoal] (f) at (8.4,-2.6) {Earlier access to\\effective therapy};
\node[mmgoal] (g) at (4.2,-2.6) {Extended PFS / OS\\(5-year OS < 13\%)};
\node[mmctx]  (n) at (8.4,-4.8) {Does not change tumor biology\\or fixed review clocks};
\draw[mmedge] (a) -- (b);
\draw[mmedge] (b) -- (c);
\draw[mmedge] (c) -- (d);
\draw[mmedge] (d) -- (e);
\draw[mmedge] (e) -- (f);
\draw[mmedge] (f) -- (g);
\draw[mmedged] (n.west) to[out=180,in=-90,looseness=0.7] (g.south);
\end{mermaidfig}
\figcaption{Figure 3. The patient time-saved cascade. Faster validated authoring
advances each step from IND submission to first dose, then through faster 3+3 cohort
turnover to the recommended Phase 2 dose, shortening the white space between phases
and, for this low-survival population (five-year overall survival under 13 percent),
extending progression-free and overall survival, without changing tumor biology or
the fixed regulatory review clocks.}

\subsection{Indication to be Studied}

The indication to be studied is histologically or cytologically confirmed pancreatic
ductal adenocarcinoma that is resectable or borderline-resectable on
multidisciplinary radiologic review, harbors a documented KRAS G12 mutation (for
example G12D, G12V, or G12R) on a validated tumor genotyping assay, occurs in an
adult of at least 18 years with an Eastern Cooperative Oncology Group (ECOG)
performance status of 0 or 1, and is anatomically compatible with the eight-arm
robotic system and with a planned pancreaticoduodenectomy. Eligibility for
perioperative daraxonrasib follows the broad pan-KRAS criteria of the RASolute 302
program \cite{daraxwhipple}. The indication is deliberately narrow at this
first-in-human stage so that the safety and feasibility readouts are interpretable:
distant metastatic disease, vascular involvement precluding a margin-negative
resection, prior pancreatic resection precluding the planned reconstruction,
uncontrolled intercurrent illness, pregnancy or lactation, and known
contraindication to daraxonrasib are exclusionary. Table 1 summarizes the
eligibility frame that defines the indication; the full inclusion and exclusion
criteria are set out in the protocol and reproduced in the proposed clinical
research module of this IND.

\noindent\begin{tabularx}{\textwidth}{L{3.0cm} Y Y}
\toprule
\textbf{Domain} & \textbf{Inclusion (must meet all)} & \textbf{Exclusion (any one disqualifies)} \\
\midrule
Disease and stage & Histologically or cytologically confirmed PDAC; resectable or borderline-resectable on multidisciplinary review. & Distant metastatic disease on staging; vascular involvement precluding R0 resection or anatomy incompatible with the eight-arm system. \\
\addlinespace
Genotype & Documented KRAS G12 mutation (G12D, G12V, or G12R) on a validated assay; eligible under pan-KRAS RASolute 302 criteria. & Absence of a qualifying KRAS G12 mutation. \\
\addlinespace
Performance and organ function & ECOG performance status 0 or 1; adequate hematologic, hepatic, and renal function for major pancreaticobiliary surgery and daraxonrasib. & Uncontrolled cardiac, pulmonary, hepatic, renal, or active infectious disease making surgery or dosing unsafe. \\
\addlinespace
Surgical candidacy & Candidate for pancreaticoduodenectomy; tumor and vascular anatomy compatible with port placement and the working volume. & Prior pancreatic resection or prior surgery precluding the planned reconstruction. \\
\addlinespace
Age and consent & Age at least 18 years; signed informed consent including explicit election or declination of the Physical AI opt-out. & Inability to consent or to comply with the Schedule of Activities. \\
\addlinespace
Other & Measured baseline CA 19-9 recorded for serial assessment (no threshold imposed). & Pregnancy or lactation; known contraindication or hypersensitivity to daraxonrasib or a required excipient; unmanageable drug interaction. \\
\bottomrule
\end{tabularx}
\figcaption{Table 1. Eligibility summary defining the indication to be studied.
Selection is equitable and is determined solely by the scientific and safety
criteria; limited English proficiency is not an exclusion, and qualified
interpreters and translated consent materials are provided.}

A synthetic reference exemplar, designated PAT-PDAC-0001, illustrates a
representative eligible participant and is used for design illustration only: a
68-year-old with a 3.4 cm head-of-pancreas tumor abutting the superior mesenteric
vein at 75 degrees, KRAS G12D, microsatellite stable, a baseline CA 19-9 of 412
U/mL, ECOG performance status 1, after a completed neoadjuvant modified FOLFIRINOX
course of four cycles. This exemplar does not represent an enrolled individual but
fixes the anatomic and oncologic context in which the eligibility frame and the
device no-fly gating are exercised \cite{daraxwhipple}.

\subsection{General Approach for Evaluation of Treatment}

The general approach is a prospective, open-label, single-arm, first-in-human study
that combines a standard Phase 1 3+3 dose-finding evaluation of perioperative
daraxonrasib with an early-feasibility evaluation of the LLM-directed eight-arm
robotic Whipple system. A single-arm, open-label design is scientifically and
ethically appropriate, and a randomized comparator arm is not, because the
first-in-human objectives, characterizing early safety, identifying a recommended
Phase 2 dose, and demonstrating that a novel surgical system can perform its
intended function within hard cyber-physical limits, are not comparative-efficacy
questions and are answered most safely in a small, intensively monitored cohort with
sentinel staggered enrollment and case-by-case data and safety monitoring board
(DSMB) review. The early-feasibility pathway recognized by FDA for significant-risk
devices is explicitly structured around small first-in-human cohorts that generate
the device-performance and safety evidence needed before any controlled comparison
\cite{phase1guidance}. The 2025 Dutch nationwide cohort of 1000 robotic
pancreaticoduodenectomies provides a robust external human benchmark for the
secondary surgical-quality endpoints, so the absence of an internal control does not
leave the surgical results uninterpretable \cite{daraxwhipple}.

\subsubsection{Objectives and Endpoints}

The study carries three co-primary objectives, one for each axis of the combined
IND/IDE investigation, beneath which sit secondary oncologic-surgical quality
endpoints benchmarked against the Dutch 2025 cohort and exploratory non-confirmatory
endpoints. The first co-primary objective characterizes the early safety of the
combined intervention; its endpoint is the incidence of device-related or
procedure-related serious adverse events through 30 days after the procedure,
including Clavien-Dindo grade III or higher complications, with device or procedure
attribution adjudicated by an independent reviewer blinded to the daraxonrasib dose
level. The second co-primary objective determines the maximum tolerated dose (MTD)
and recommended Phase 2 dose (RP2D) of perioperative daraxonrasib, defined
operationally by the rate of dose-limiting toxicities (DLTs) within the 28-day DLT
window of the 3+3 escalation. The third co-primary objective establishes the
technical feasibility of the eight-arm robotic Whipple under continuous human
oversight; its endpoint is the proportion of procedures that complete all
prespecified assigned operative tasks without unplanned conversion caused by device
malfunction or unsafe device behavior. Figure 4 renders the objective-to-endpoint
hierarchy, and Table 2 sets out the full mapping with justifications and the Dutch
2025 comparators \cite{daraxwhipple,phase1guidance}.

\begin{mermaidfig}
\node[mmgoal] (pr) at (0,0)       {PRIMARY (co-primary)\\safety + feasibility};
\node[mmstep] (p1) at (-4.6,-2.6) {Device / procedure SAE\\incidence through 30 days};
\node[mmstep] (p2) at (0,-2.6)    {MTD / RP2D of\\daraxonrasib\\(3+3, 28-day DLT)};
\node[mmstep] (p3) at (4.6,-2.6)  {Assigned tasks completed\\without unsafe\\conversion};
\node[mmaccent](se)at (0,-5.2)    {SECONDARY\\oncologic-surgical};
\node[mmin]   (s1) at (-4.6,-7.6) {R0 rate; ISGPS\\B/C fistula};
\node[mmin]   (s2) at (0,-7.6)    {Clavien-Dindo III+;\\30 / 90-day mortality};
\node[mmin]   (s3) at (4.6,-7.6)  {Conversion rate; operative\\parameters (vs Dutch 2025)};
\node[mmctx]  (ex) at (0,-10.0)   {EXPLORATORY\\non-confirmatory};
\node[mmin]   (e1) at (0,-12.0)   {PFS / OS to 24 mo (RECIST);\\LLM concordance;\\sim-to-real gap};
\draw[mmedge] (pr) -- (p1);
\draw[mmedge] (pr) -- (p2);
\draw[mmedge] (pr) -- (p3);
\draw[mmedged] (pr) -- (se);
\draw[mmedge] (se) -- (s1);
\draw[mmedge] (se) -- (s2);
\draw[mmedge] (se) -- (s3);
\draw[mmedged] (se) -- (ex);
\draw[mmedge] (ex) -- (e1);
\end{mermaidfig}
\figcaption{Figure 4. The objective-to-endpoint hierarchy of the general
investigational plan. Three co-primary objectives (device or procedure
serious-adverse-event incidence through 30 days, the maximum tolerated dose and
recommended Phase 2 dose of daraxonrasib by 3+3 over the 28-day window, and
completion of assigned operative tasks without unsafe conversion) sit above
secondary oncologic-surgical quality endpoints (R0 rate, ISGPS grade B/C fistula,
Clavien-Dindo grade III or higher, 30 and 90-day mortality, conversion rate, and
operative parameters benchmarked against the Dutch 2025 robotic cohort) and
exploratory, non-confirmatory endpoints (progression-free and overall survival to 24
months by RECIST, LLM advisory concordance, and the sim-to-real gap).}

\noindent\begin{tabularx}{\textwidth}{L{2.4cm} L{4.7cm} Y}
\toprule
\textbf{Tier} & \textbf{Objective and endpoint} & \textbf{Justification / comparator} \\
\midrule
\multicolumn{3}{l}{\textit{Co-primary}} \\
\midrule
Safety & Device-related or procedure-related SAE incidence through 30 days, including Clavien-Dindo grade III+, dose-blinded attribution. & Standard first-in-human safety anchor; captures the early lethal cascade of fistula, hemorrhage, sepsis, and failure to rescue. \\
\addlinespace
Dose-finding & MTD and RP2D of perioperative daraxonrasib by the DLT rate within the 28-day window of the 3+3 across DL1 to DL3. & Established dose-finding design for targeted oncology agents; anchored to the RASolute 302 pan-KRAS program. \\
\addlinespace
Feasibility & Proportion of procedures completing all assigned operative tasks without unplanned conversion from device malfunction or unsafe behavior. & FDA-recognized early-feasibility readout for a significant-risk device. \\
\midrule
\multicolumn{3}{l}{\textit{Secondary (descriptive, Dutch 2025 comparators)}} \\
\midrule
Oncologic adequacy & R0 (microscopically margin-negative) resection rate at day-90 pathology. & Strongest modifiable surgical determinant of survival; descriptive comparison with Dutch 2025. \\
\addlinespace
Fistula & ISGPS grade B/C postoperative pancreatic fistula rate. & Leading driver of post-Whipple morbidity; Dutch 2025 comparator 24.4 percent. \\
\addlinespace
Major complications & Clavien-Dindo grade III+ complication rate through 90 days. & Validated severity threshold; Dutch 2025 comparator 41.3 percent. \\
\addlinespace
Mortality & 30-day and 90-day all-cause mortality. & Standard short-term survival safety measures; Dutch 2025 30-day comparator 3.9 percent. \\
\addlinespace
Conversion & Unplanned conversion rate from the robotic system to open or conventional approach. & Contextualizes feasibility; Dutch 2025 comparator 10.1 percent. \\
\addlinespace
Operative parameters & Estimated blood loss, operative time, and length of stay. & Mean and median with range; Dutch 2025 learning-curve context. \\
\midrule
\multicolumn{3}{l}{\textit{Exploratory (non-confirmatory)}} \\
\midrule
Survival & PFS and OS assessed to 24 months by RECIST v1.1, Kaplan-Meier. & Preliminary, hypothesis-generating signal; the small single-arm sample precludes confirmatory inference. \\
\addlinespace
Advisory reliability & Concordance between the LLM advisory output and the independent hard-coded safety gate. & Characterizes the second-opinion oracle and informs autonomy-graduation decisions reserved for later phases. \\
\addlinespace
Simulation fidelity & Sim-to-real gap in trajectory (mm) and tip force (N). & Validates the Phase 0 simulation-validation gate; targets less than 2 mm and less than 0.5 N. \\
\bottomrule
\end{tabularx}
\figcaption{Table 2. The objective-to-endpoint mapping with justifications and the
Dutch 2025 nationwide robotic-pancreatoduodenectomy comparators. Comparisons are
descriptive only; the single-arm design supports estimation, not inferential
between-group testing.}

\subsubsection{Dose-Escalation Approach}

Daraxonrasib is administered orally once daily. The escalation uses exact doses
anchored to the RASolute 302 pan-KRAS program, in which the 300 mg once-daily dose
is the established RP2D. To preserve a conservative first-in-human margin in the
perioperative setting, the escalation begins below that reference dose and steps up
to it: DL1 is 160 mg once daily (the starting dose), DL2 is 220 mg once daily, and
DL3 is the 300 mg once-daily reference dose. Three participants enroll at each level
and are observed over a 28-day DLT window. If no DLT occurs among the first three,
the cohort escalates; if one of three experiences a DLT, the cohort expands to six;
a single DLT among six permits escalation, whereas two or more DLTs at a level
define that level as exceeding the MTD and fix the MTD at the preceding level. The
MTD is the highest level at which fewer than two of six DLT-evaluable participants
experience a DLT, and the RP2D is selected at or below the MTD with integration of
all safety, tolerability, and available pharmacokinetic data. Figure 5 renders the
decision automaton and Table 3 states the rule explicitly for each cohort outcome
\cite{daraxwhipple,phase1guidance}.

\begin{mermaidfig}
\node[mmin]  (dl1) at (0,0)       {DL1: 160 mg QD\\enroll 3};
\node[mmdec] (g1)  at (0,-2.6)    {DLT in\\28-day window?};
\node[mmstep](e1)  at (4.6,-2.6)  {Expand to 6\\at this level};
\node[mmin]  (dl2) at (0,-5.4)    {DL2: 220 mg QD\\enroll 3};
\node[mmdec] (g2)  at (0,-8.0)    {DLT?};
\node[mmstep](e2)  at (4.6,-8.0)  {Expand to 6};
\node[mmin]  (dl3) at (0,-10.8)   {DL3: 300 mg QD\\enroll 3};
\node[mmdec] (g3)  at (0,-13.4)   {DLT?};
\node[mmstep](e3)  at (4.6,-13.4) {Expand to 6};
\node[mmdark](mtd) at (9.4,-8.0)  {MTD = prior level\\($\geq$ 2 of 6 DLT)};
\node[mmgoal](rp)  at (9.4,-13.4) {RP2D at or below MTD\\integrating PK + safety};
\draw[mmedge] (dl1) -- (g1);
\draw[mmedge] (g1) -- node[left,font=\scriptsize]{0 of 3} (dl2);
\draw[mmedge] (g1) -- node[above,font=\scriptsize]{1 of 3} (e1);
\draw[mmedge] (e1) to[out=-90,in=0,looseness=0.6] node[right,font=\scriptsize,pos=0.2]{1 of 6} (dl2.east);
\draw[mmedge] (e1) to[out=0,in=90,looseness=0.7] node[right,font=\scriptsize,pos=0.3]{$\geq$ 2 of 6} (mtd.north);
\draw[mmedge] (dl2) -- (g2);
\draw[mmedge] (g2) -- node[left,font=\scriptsize]{0 of 3} (dl3);
\draw[mmedge] (g2) -- node[above,font=\scriptsize]{1 of 3} (e2);
\draw[mmedge] (e2) to[out=-90,in=180,looseness=0.6] node[left,font=\scriptsize,pos=0.3]{1 of 6} (dl3.west);
\draw[mmedge] (e2) -- node[above,font=\scriptsize]{$\geq$ 2 of 6} (mtd);
\draw[mmedge] (dl3) -- (g3);
\draw[mmedge] (g3) -- node[above,font=\scriptsize]{0 of 3 / 1 of 6} (rp);
\draw[mmedge] (e3) to[out=90,in=-90,looseness=0.7] node[right,font=\scriptsize,pos=0.4]{$\geq$ 2 of 6} (mtd.south);
\draw[mmedge] (g3) -- node[above,font=\scriptsize,pos=0.6]{1 of 3} (e3);
\draw[mmedge] (mtd) -- (rp);
\end{mermaidfig}
\figcaption{Figure 5. The standard 3+3 dose-escalation automaton across the three
daraxonrasib dose levels (DL1 160 mg, DL2 220 mg, DL3 300 mg once daily). Three
participants enroll at each level; with no dose-limiting toxicity in the 28-day
window the cohort escalates, with one of three it expands to six, and two or more of
six fixes the maximum tolerated dose at the prior level. The recommended Phase 2 dose
is selected at or below the maximum tolerated dose, integrating pharmacokinetic and
safety data. The design maximum is three levels times six, up to 18 treated.}

\noindent\begin{tabularx}{\textwidth}{L{2.6cm} L{2.4cm} L{3.4cm} Y}
\toprule
\textbf{Cohort observation} & \textbf{DLT count} & \textbf{Decision} & \textbf{Consequence} \\
\midrule
First 3 at a level & 0 of 3 & Escalate & Open the next higher dose level; current level deemed tolerated at the cohort of 3. \\
\addlinespace
First 3 at a level & 1 of 3 & Expand to 6 & Enroll 3 additional participants at the same level before any escalation decision. \\
\addlinespace
First 3 at a level & $\geq$ 2 of 3 & MTD exceeded & Do not enroll further at this level; the MTD is the preceding (lower) level. \\
\addlinespace
Expanded 6 at a level & 1 of 6 & Escalate & Single DLT in six permits escalation to the next higher level. \\
\addlinespace
Expanded 6 at a level & $\geq$ 2 of 6 & MTD exceeded & The level exceeds the MTD; the MTD is the preceding level. \\
\addlinespace
Highest level (DL3) & 0 of 3 or 1 of 6 & RP2D candidate & DL3 (300 mg) tolerated; RP2D selected at or below DL3 integrating PK and safety. \\
\addlinespace
Any level & $\geq$ 2 of 6 & Declare MTD / RP2D & MTD fixed at prior level; RP2D at or below MTD. \\
\bottomrule
\end{tabularx}
\figcaption{Table 3. The 3+3 dose-escalation decision rule. The maximum tolerated
dose is the highest level at which fewer than two of six DLT-evaluable participants
experience a dose-limiting toxicity within the 28-day window; the recommended Phase
2 dose is selected at or below the MTD. The design maximum is three levels times six,
up to 18 treated.}

\subsubsection{Device Readiness and Statistical Framework}

The device has no pharmacologic dose; its dose analogue is the readiness evidence
required before any patient-facing robotic activity, escalated only by passing a
mandatory Phase 0 simulation-validation gate. Three conditions must be met and
documented before the first robotic case: a Unified Safety Level (USL) rating of at
least 7.0 on the surgical-readiness scale; completion of at least 1000 simulated
procedures across at least two independent frameworks (for example Isaac, MuJoCo,
Gazebo, or PyBullet), so the finding is not an artifact of any single simulator; and
a sim-to-real fidelity within prespecified limits, a trajectory gap less than 2 mm
and a tip-force gap less than 0.5 N between the simulation prediction and the
operative record. The gate is re-validated after any change to the system. These
limits operate together with the intra-operative hard caps (a cross-arm
emergency-stop budget at or below 3 ms, per-arm tip-force caps at or below 3 N, a
cumulative cross-arm cap at or below 18 N, and vascular no-fly gating) so the device
dose analogue is met only when both the simulation evidence and the real-time safety
envelope are demonstrably in place \cite{onpremwhippl,daraxwhipple}.

The statistical framework is that of a small, intensively monitored, single-arm
investigation whose primary purpose is to characterize early safety and to identify
a recommended dose, not to test a comparative-efficacy hypothesis. The analyses are
predominantly descriptive and estimation-based, follow ICH E9, and are reported in
accordance with TRIPOD+AI for the model-assisted advisory component
\cite{ichpaistduni}. No power calculation is performed because the study tests no
confirmatory efficacy hypothesis. Every proportion is accompanied by an exact
(Clopper-Pearson) two-sided 95 percent confidence interval rather than a
normal-approximation interval, which is unreliable near the bounds and at small
sample sizes; time-to-event variables are summarized by Kaplan-Meier with the number
at risk; missing data are not imputed for the primary or safety analyses; and a
nominal alpha of 0.05 two-sided applies to supportive estimation only, with no
multiplicity adjustment because the analyses are descriptive and non-confirmatory.
Four analysis populations are applied consistently: the Safety population (all who
received any intervention), the DLT-evaluable population (all dosed participants who
completed the 28-day window or had a DLT within it), the Per-Protocol population,
and the modified intention-to-treat population \cite{ichpaistduni,phase1guidance}.

\subsection{Description of First Year Trial(s)}

During the first year, a single trial is conducted at a single academic medical
center: the combined IND/IDE first-in-human study of perioperative daraxonrasib and
the LLM-directed eight-arm robotic Whipple. Enrollment is planned for up to 18
treated participants, allocated across the three daraxonrasib dose levels (DL1
through DL3) by the 3+3 rule with staggered sentinel enrollment, from approximately
36 screened. Recruitment proceeds at approximately one to two participants per month
through the institutional multidisciplinary pancreatic-cancer program, with eligible
individuals identified at tumor-board review and described to them by a member of the
research team who is not the treating surgeon, to reduce undue influence. The planned
accrual rate, the sentinel staggered enrollment of the 3+3 design, and the Phase 0
readiness gate together give an enrollment period sufficient to treat 18 participants
from approximately 36 screened \cite{daraxwhipple,phase1guidance}.

Sentinel staggering is a binding feature of the first-year conduct: no two
participants in a cohort undergo surgery within the same DLT-evaluation interval
until the supervising surgeon and the DSMB have reviewed the preceding case. After
each participant completes the 28-day DLT-evaluation window and before the next
participant in the cohort undergoes surgery, the accumulating safety data are
reviewed; after each dose-level cohort is complete, the DSMB convenes for a formal
interim safety review and an escalation, expansion, hold, or halt decision. The
prespecified halt rules are quantitative and binding: enrollment and escalation halt
for a mandatory DSMB review if the device-related or procedure-related serious-AE
rate reaches or exceeds one third within a cohort, or if two device-related deaths
occur at any point. Additional pause triggers include any unexpected device-related
serious AE, any breach of a hard cyber-physical safety limit (emergency-stop
latency, positional, or force envelope), and any failure of the pre-procedure
safety-matrix verification. The DSMB may halt, pause, modify, or permit continuation,
and a halt for an immediate hazard is implemented at once \cite{daraxwhipple}.

The device operates within a phase-graduated staged-autonomy model. In this Phase 1
study the system runs only at the clinician-controlled (teleoperation or supervised)
level and the supervised semiautonomous level, both under continuous human oversight
with an immediate manual override; full autonomy is prohibited at this phase
consistent with 21 CFR \S 312.21(e). Any future move to higher autonomy (Stage 2
supervised semiautonomous expansion or Stage 3 higher autonomy) is reserved for a
later phase and gated by independent safety review, a USL reassessment, and DSMB
concurrence. The perioperative daraxonrasib pause-and-restart advisory is exercised
in every case: daraxonrasib is paused before surgery and the restart day is set by
an LLM-bound, human-confirmed advisory whose inputs are the postoperative pancreatic
fistula grade by ISGPS (A, B, or C) and the serum daraxonrasib trough relative to a
0.5 ng/mL threshold, with outputs of T+7 days, T+14 days, T+21 days, or hold. A
32-iteration deterministic sweep returned a baseline serum of 0.45 ng/mL
(sub-threshold) in all iterations, a T+7-day restart in 29 of 32 (90.6 percent,
Grade A and sub-threshold), a T+14-day restart in 3 of 32 (9.4 percent, Grade B or
delayed), and no T+21-day restart \cite{daraxwhipple,clintrustpai}.

The first-year objectives map onto the endpoint hierarchy of Figure 4: the
co-primary safety, dose-finding, and feasibility readouts are accrued as cohorts
complete, the secondary surgical-quality endpoints are estimated against the Dutch
2025 benchmark as the per-protocol set matures, and the exploratory survival,
advisory-concordance, and sim-to-real measures begin to accumulate through the day-90
pathology review and the 24-month survival follow-up. The realistic first-year
deliverable is completion of the DL1 and DL2 cohorts and initiation of the DL3
cohort, with an interim determination of the MTD and a provisional RP2D if DL3 is
reached and tolerated, together with the device-feasibility and telemetry summaries
that support progression to a later-phase efficacy study \cite{daraxwhipple}.

\subsection{Number of Subjects to be Evaluated}

The planned sample size is up to 18 treated participants, allocated across three
daraxonrasib dose levels (DL1 through DL3) by the 3+3 rule, from approximately 36
screened. This number is not derived from a power calculation, because the study
tests no confirmatory efficacy hypothesis; it is determined jointly by the
dose-finding design and by the early-feasibility precedent for significant-risk
devices, and it is justified by its operating characteristics. The 3+3 rule enrolls
cohorts of three: a level with no DLT in three escalates, a level with one DLT in
three expands to six, and a level with two or more DLTs in up to six is exceeded,
defining the dose below it as the MTD. The maximum exposure at any single dose level
is therefore six participants, and three escalation levels give a design maximum of
three times six, equal to 18 treated participants; many realizations terminate with
fewer. This ceiling is consistent with the FDA early-feasibility convention of
generally fewer than ten participants per device-feasibility cohort while providing
the per-dose exposure the drug-arm 3+3 escalation requires \cite{phase1guidance}.

Rather than power, the design is characterized by its operating characteristics and
by the precision of estimation at the planned sample sizes. The probability that the
3+3 rule escalates past a given dose level is a function of the true DLT probability
at that level: a dose with a true DLT rate near 10 percent is very likely to be
tolerated and escalated, whereas a dose with a true DLT rate of 40 percent or higher
is very likely to be identified as excessively toxic and to define the MTD at the
level below. The exact-interval precision is summarized in Table 4. Zero events in
three yields an exact two-sided 95 percent confidence interval of approximately 0 to
71 percent; zero in six, approximately 0 to 46 percent; zero in 18, approximately 0
to 18 percent; and one event in 18, a point estimate of 5.6 percent with an interval
of approximately 0.1 to 27 percent. These intervals make explicit both the value and
the limits of a first-in-human cohort: a clean safety record across the full cohort
of 18 bounds the true event rate below roughly one in five, while a single event
still excludes rates above roughly one quarter \cite{phase1guidance,ichpaistduni}.

\noindent\begin{tabularx}{\textwidth}{L{2.8cm} L{3.0cm} Y}
\toprule
\textbf{Quantity} & \textbf{Planned value} & \textbf{Operating characteristic / precision} \\
\midrule
Cohort size & 3, expand to 6 & 3+3 rule: escalate on 0/3; expand on 1/3; exceed dose on $\geq$ 2/6. \\
\addlinespace
Per-dose maximum & 6 & Maximum exposure at any one dose level before the MTD is declared. \\
\addlinespace
Dose levels & 3 (DL1 to DL3) & Three-level escalation; design maximum 3 x 6 = 18 treated. \\
\addlinespace
Total treated & up to 18 (approx. 36 screened) & At or below the early-feasibility convention of generally fewer than 10 per device cohort, summed across dose levels. \\
\addlinespace
0 events / 3 & point estimate 0\% & Exact two-sided 95\% CI approximately 0\% to 71\%. \\
\addlinespace
0 events / 6 & point estimate 0\% & Exact two-sided 95\% CI approximately 0\% to 46\%. \\
\addlinespace
0 events / 18 & point estimate 0\% & Exact two-sided 95\% CI approximately 0\% to 18\%. \\
\addlinespace
1 event / 18 & point estimate 5.6\% & Exact two-sided 95\% CI approximately 0.1\% to 27\%. \\
\addlinespace
Significance level & $\alpha$ = 0.05 (two-sided) & Applies to supportive estimation only; no $\alpha$ is spent on dose selection. \\
\bottomrule
\end{tabularx}
\figcaption{Table 4. Operating characteristics and exact-interval (Clopper-Pearson)
precision at the planned sample sizes. The dose-finding behavior is that of the
standard 3+3 rule; no statistical power is computed because the study tests no
confirmatory efficacy hypothesis.}

The four analysis populations apportion the 18 participants across the analyses
without altering the total exposure: the Safety population comprises all who received
any intervention and is the primary population for safety and primary-endpoint
analyses; the DLT-evaluable population comprises those who completed the 28-day
window or had a DLT within it and drives the escalation and MTD determination; the
Per-Protocol population supports the secondary surgical-quality analyses; and the
modified intention-to-treat population supports the supportive and exploratory
efficacy summaries. Because the cohort is small and intensively monitored,
individual participant data are tabulated in full, so that every aggregate statistic
can be traced to its contributing records \cite{ichpaistduni}.

\subsection{Drug Related Risks}

The foreseeable risks of the investigation arise from three sources, the
investigational drug, the investigational device, and the surgical procedure, and
each is bounded by a specific control rather than left to clinical judgment alone.
The anticipated risks of daraxonrasib are those characteristic of RAS-pathway
inhibition and of the RASolute 302 experience, including gastrointestinal effects
(nausea, diarrhea, and stomatitis), dermatologic effects (rash), hepatic enzyme
elevation, and the class effects of MAPK-pathway suppression; in the perioperative
setting the additional concern is impaired wound healing and anastomotic integrity,
which is the explicit reason for the pause-and-restart advisory that holds
daraxonrasib before surgery and resumes it only on a human-confirmed,
fistula-grade-and-trough-informed schedule (T+7, T+14, T+21 days, or hold). Drug
toxicity is captured through the 28-day DLT window, the dose-blinded SAE attribution,
the exact-interval safety reporting, and the binding cohort-level halt rules
\cite{daraxwhipple,oncotrialllm}. Table 5 summarizes the principal risks, their
sources, and their mitigations.

\noindent\begin{tabularx}{\textwidth}{L{3.4cm} Y Y}
\toprule
\textbf{Risk source} & \textbf{Foreseeable risk} & \textbf{Mitigation} \\
\midrule
Daraxonrasib (drug) & GI effects (nausea, diarrhea, stomatitis), rash, hepatic enzyme elevation, MAPK-pathway class effects; impaired perioperative wound healing. & 28-day DLT window; conservative 160 mg starting dose below the 300 mg RP2D; pause-and-restart advisory keyed to ISGPS grade and 0.5 ng/mL trough; cohort-level halt rules. \\
\addlinespace
Robotic system (device) & Device malfunction, unsafe trajectory, excessive tip force, vascular injury near the SMV, PV, HA, CA, or SMA. & Force caps $\leq$ 3 N per arm and $\leq$ 18 N cumulative; cross-arm E-stop $\leq$ 3 ms; hardware E-stop $\leq$ 500 ms; 10 kHz heartbeat with 100 microsecond watchdog; vascular no-fly gating; Phase 0 gate (USL $\geq$ 7.0, $\geq$ 1000 sims, sim-to-real $<$ 2 mm and $<$ 0.5 N). \\
\addlinespace
LLM advisory layer & Erroneous or unverified advisory output; over-reliance (therapeutic misconception); model degradation. & Second-opinion advisory only; full autonomy prohibited (\S 312.21(e)); ten-gate verification-before-generation (ACCEPT/BLOCK/ESCALATE); hash-pinned deterministic seed and commit; Physical AI opt-out; advisory-to-gate concordance monitored. \\
\addlinespace
Whipple procedure & Pancreatic fistula, hemorrhage, intra-abdominal infection, sepsis, failure to rescue, unplanned conversion. & Sentinel staggered enrollment with case-by-case DSMB review; three-tier oversight (DSMB, Independent Safety Monitor, Physical AI Safety Review Committee); conversion pathway with continuous human override. \\
\addlinespace
Schedule and population & Delay to effective therapy in a low-survival disease. & Verified single-prompt authoring compresses only the administrative bucket, advancing first dose without altering tumor biology or fixed review clocks (Figure 3). \\
\bottomrule
\end{tabularx}
\figcaption{Table 5. Principal drug-related, device-related, procedural, and
schedule-related risks with their controls. Each risk is bounded by a specific,
prespecified mitigation; the residual risk after these controls is judged acceptable
for the population and the stage of investigation.}

The device-related and procedural risks are bounded by the cyber-physical safety
envelope and the three-tier oversight structure. The hard caps (per-arm tip force at
or below 3 N, cumulative cross-arm force at or below 18 N, cross-arm emergency stop
within 3 ms to 0.0 N, and hardware emergency stop within 500 ms), the 10 kHz
heartbeat with the 100 microsecond watchdog, and the five-vessel no-fly gating limit
the magnitude of any single excursion, while the Phase 0 gate (USL at least 7.0, at
least 1000 simulations across at least two frameworks, and a sim-to-real gap below 2
mm and 0.5 N) ensures the system is demonstrably ready before any patient-facing
activity. The LLM advisory risks (an erroneous or unverified advisory and the
therapeutic misconception of over-reliance on the system) are bounded by the
advisory-only role, the prohibition on full autonomy under 21 CFR \S 312.21(e), the
ten-gate verification-before-generation suite, the hash-pinned deterministic seed
and commit, the Physical AI opt-out in the consent, and the monitoring of
advisory-to-gate concordance. Safety reporting follows the 21 CFR \S 312.32 clocks
(fatal or life-threatening suspected adverse reactions within 7 calendar days and
other serious and unexpected reactions within 15 calendar days), augmented by the
six Physical AI reporting triggers under \S 312.32(g) and an audit trail preserved
from 24 hours before to 72 hours after each Physical AI adverse event under 21 CFR
part 11 \cite{cfr312paiuni,congress9510}.

The benefit-risk balance is favorable for this population and at this stage. The
population is vulnerable, with a lethal malignancy whose five-year survival is below
13 percent and for whom the resection margin is the strongest modifiable determinant
of survival; the intervention pairs a genotype-matched systemic agent with a
margin-negative curative-intent resection delivered within hard cyber-physical limits
under continuous human oversight. The schedule benefit reaches the patient through
the cascade of Figure 3: faster validated authoring advances IND and IRB submission
and site activation, the elapse of the 30-day clock and the first dose, and the 3+3
cohort turnover toward the RP2D, shortening the white space between phases and, for
this low-survival population, extending progression-free and overall survival,
without changing tumor biology or any fixed review clock. For patients whose disease
will not pause, the alternative of slower sequential authoring carries its own,
often larger, cost in delayed treatment. With each foreseeable risk bounded by a
prespecified control, with the conservative 160 mg starting dose, the sentinel
staggered enrollment, the binding halt rules, and the three-tier oversight, the
probable benefit of the investigation exceeds its probable risk, and the General
Investigational Plan supports the proposed first-year conduct \cite{pdac060s2030,daraxwhipple,clintrustpai}.

\subsection{References}

The clinical, statistical, and acceleration content of this plan is grounded in the
trial protocol and the supporting evidence base. The Phase 1 first-in-human guidance
and the early-feasibility framework inform the single-arm dose-finding design and the
operating characteristics \cite{phase1guidance}. The daraxonrasib robotic Whipple
program and the on-premises eight-arm system specifications supply the drug, device,
and advisory facts \cite{daraxwhipple,onpremwhippl,oncotrialllm}. The Physical AI
adaptations of 21 CFR Part 312 and Part 50 and the ICH guidance govern the safety
reporting, consent, and statistical conventions \cite{cfr312paiuni,cfr050paiuni,ichpaistduni}.
H.R. 9510 and the clinical-trust framework establish the verification-before-generation
discipline \cite{congress9510,clintrustpai}. The PDAC epidemiology, the acceleration
mechanism, and the supporting computational evidence are drawn from the population
and document-acceleration sources \cite{pdac060s2030,paipredict4x,natlpaiplatf,paisitedocpk,paiautospons,llmcomp16fda,fdadigtwinpc,qspmetpancre,chatgpt100kp,pdacdigtwinp}.
The full bibliography is consolidated in the references and back-matter section of
this IND.
